\section{トルク抑制のプロセス}
\label{sec:method_overview}

\subsection{LIPM に基づく CoM 高さと横方向転倒挙動の関係}
\label{subsec:lipm_height_effect}

二足歩行ロボットの DS から SS への遷移時に生じる転倒挙動は，特にロール方向，すなわち左右方向の重心運動として顕在化する．ここではまず，線形倒立振子モデル（Linear Inverted Pendulum Model: LIPM）の基本式に立ち返り，重心高さが横方向の重心加速度に与える影響について整理する．

LIPM では，重心高さを一定値 $z_{\mathrm{c}}$ と仮定し，角運動量の影響を無視することで，CoM と ZMP の関係は次式で表される（ \cref{eq:lipm} の $y$ 軸方向）．
\begin{equation}
\ddot{y} = \frac{g}{z_{\mathrm{c}}}\left(y - y_{\mathrm{zmp}}\right)
\label{eq:lipm_y}
\end{equation}
ここで，$y$ は CoM の左右方向位置，$y_{\mathrm{zmp}}$ は ZMP の左右方向位置，$g$ は重力加速度である．式 \eqref{eq:lipm_y} は，ZMP からの相対変位 $\left(y - y_{\mathrm{zmp}}\right)$ に比例して，横方向の重心加速度が生じることを示している．

式 \eqref{eq:lipm_y} の比例係数は $g / z_{\mathrm{c}}$ であり，これは重心高さ $z_{\mathrm{c}}$ に反比例する．したがって，$z_{\mathrm{c}}$ が増大すると，同一の ZMP--CoM 偏差に対して生じる加速度 $\ddot{y}$ は小さくなる．すなわち，重心高さを高くすることは，横方向の重心運動の加速を抑制し，転倒方向への倒れ込みを緩やかにする方向に働く．

\subsection{振り子モデルに基づく CoM 高さと転倒挙動の力学的解釈}
\label{subsec:pendulum_theory}

片脚支持期における二足歩行ロボットのロール方向挙動は，LIPMに基づき，支持脚足先を支点とした倒立振り子運動として近似的に捉えることができる．ここでは，CoM を質点とし，支点から CoM までの距離を $\ell$ とする単純な振り子モデルを用いて，CoM 高さが転倒挙動に与える影響を力学的に考察する．

ロール方向の傾き角を $\theta$ とすると，支点まわりの運動方程式は次式で表される．
\begin{equation}
I \ddot{\theta} = m g \ell \sin\theta
\end{equation}
ここで，$m$ はロボットの質量，$g$ は重力加速度，$I$ は支点まわりの慣性モーメントである．質点近似を用いると，慣性モーメントは
\begin{equation}
I = m \ell^2
\end{equation}
と表されるため，
\begin{equation}
m \ell^2 \ddot{\theta} = m g \ell \sin\theta
\end{equation}
を得る．

ロール角が小さい範囲では $\sin\theta \approx \theta$ と近似でき，角加速度は
\begin{equation}
\ddot{\theta} = \frac{g}{\ell}\theta
\label{eq:angular_accel_l}
\end{equation}
と表される．式 \eqref{eq:angular_accel_l} から，角加速度 $\ddot{\theta}$ は振り子長 $\ell$ に反比例することが分かる．すなわち，CoM 高さが高くなり $\ell$ が大きくなるほど，同一の傾き角 $\theta$ に対する角加速度は小さくなり，ロボットの倒れ込み運動は緩やかになる．

この結果は，重力によるモーメントが $\ell$ に比例して増加する一方で，慣性モーメントが $\ell^2$ に比例して増加するため，振り子長が長くなるほど回転運動に対する慣性の影響が支配的となることを意味している．その結果，姿勢の変化速度が抑制され，転倒方向への傾きが急激に進行しにくくなる．

制御の観点から見ると，傾き角の時間変化が緩やかになることで，姿勢誤差や角速度の増大が抑えられ，それに対抗するために必要となる支持脚ロール方向トルクの急激な立ち上がりが緩和される．したがって，CoM 高さは，DS から SS への遷移時における転倒挙動や支持脚ロール方向トルクの大きさに強く関与する重要なパラメータであるといえる．

なお，上記の議論は，ロール方向挙動を単純化した振り子モデルに基づくものであり，二足歩行ロボットの全身多自由度ダイナミクスを厳密に表すものではない．しかしながら，LIPM に基づく横方向加速度の議論と整合的であり，CoM 高さと転倒挙動の関係を直感的かつ力学的に理解する上で，有用な近似であると考えられる．
